%%========================================================================
%%  File   : cex_contagion_v2.0/manuscript/main_eca_v2_spine.tex
%%  Purpose: Immutable spine (title + abstract + §1) baseline for v2.0 regeneration
%%  Source : Extracted from cex_contagion_v0.1/manuscript/main_eca_v01.tex (v0.9-m)
%%  Date   : 2026-08-18
%%  Note   : This file is READ-ONLY. Regeneration in P2 must reuse this spine verbatim
%%           except for character-level tweaks required by empirical consistency.
%%========================================================================

\documentclass[12pt]{article}

%---- packages -----------------------------------------------------------
\usepackage[utf8]{inputenc}
\usepackage[T1]{fontenc}
\usepackage[english]{babel}
\usepackage{amsmath,amssymb,amsthm,mathtools,mathrsfs}
\usepackage[margin=1.25in]{geometry}
\usepackage{setspace}
\onehalfspacing
\usepackage{graphicx}
\usepackage{booktabs}
\usepackage{caption}
\usepackage{float}
\usepackage{hyperref}
\usepackage[authoryear,round]{natbib}
\usepackage{titling}
\usepackage{titlesec}
\hypersetup{colorlinks=true,linkcolor=black,citecolor=black,urlcolor=black}
% \usepackage{lineno}\linenumbers  % ECA does not require line numbers

% Match sample's title layout
\pretitle{\begin{center}\LARGE\bfseries}
\posttitle{\end{center}\vspace{0.5em}}
\preauthor{\begin{center}\large}
\postauthor{\end{center}}
\predate{}
\postdate{}
\date{}

% Section headings roman-bold, no "\S" prefix
\titleformat*{\section}{\Large\bfseries}
\titleformat*{\subsection}{\large\bfseries}

%---- theorem environments (match sample: "Definition 1 (Name).") --------
\newtheoremstyle{sampledef}%
  {\topsep}{\topsep}%             margins
  {\normalfont}%                  body
  {}%                             indent
  {\bfseries}%                    head
  {.}%                            punct
  { }%                            space
  {\thmname{#1}~\thmnumber{#2}\thmnote{ (\normalfont\itshape#3\normalfont)}}
\theoremstyle{sampledef}
\newtheorem{definition}{Definition}
\newtheorem{assumption}[definition]{Assumption}

\newtheoremstyle{sampleplain}%
  {\topsep}{\topsep}%
  {\itshape}%                      body italic (theorem-style)
  {}%
  {\bfseries}%
  {.}%
  { }%
  {\thmname{#1}~\thmnumber{#2}\thmnote{ (\normalfont\itshape#3\normalfont)}}
\theoremstyle{sampleplain}
\newtheorem{theorem}[definition]{Theorem}
\newtheorem{proposition}[definition]{Proposition}
\newtheorem{lemma}[definition]{Lemma}
\newtheorem{corollary}[definition]{Corollary}

\newtheoremstyle{sampleremark}%
  {\topsep}{\topsep}%
  {\normalfont}%
  {}%
  {\itshape}%                      head italic (remark-style)
  {.}%
  { }%
  {\thmname{#1}~\thmnumber{#2}\thmnote{ (\normalfont\itshape#3\normalfont)}}
\theoremstyle{sampleremark}
\newtheorem{remark}[definition]{Remark}

%---- shortcuts ---------------------------------------------------------
\newcommand{\R}{\mathbb{R}}
\newcommand{\N}{\mathbb{N}}
\newcommand{\Z}{\mathbb{Z}}
\newcommand{\E}{\mathbb{E}}
\newcommand{\Prob}{\mathbb{P}}
\newcommand{\1}{\mathbf{1}}
\DeclareMathOperator*{\argmin}{arg\,min}

%========================================================================
\begin{document}

\title{Unlikely Intersections in\\ Crypto Exchange Reserves:\\
An O-Minimal Test for Systemic Risk}

\author{Hongjun Gou\thanks{Industrial and Commercial Bank of China,
  Beijing 100140, China. Email: \texttt{gouhongjun\_bs@cq.icbc.com.cn}.}}

\maketitle

%---- Abstract as a standalone block, sample-style ----------------------
\begin{center}\textbf{Abstract}\end{center}

Six major centralised cryptocurrency exchanges collapsed within
seven months of each other. The cluster exposed a gap in systemic
risk measurement that equity-based indices cannot fill: most crypto
venues lack liquid equity. This paper builds a supervisory
diagnostic from public reserve compositions. The joint reserve
state is embedded in a product of simplices, and simultaneous
distress corresponds to an unlikely intersection in the sense of
the o-minimality programme in arithmetic geometry. Matching upper
and lower bounds meet at a closed-form crossover threshold that
separates a common systemic factor from idiosyncratic failure. The
spot-Bitcoin ETF approval provides a sharp natural experiment. The
full intersection attains its maximum in the approval quarter,
whereas a stablecoin placebo shows no comparable jump. The
diagnostic is computable in real time and yields a state-contingent
capital buffer.

\vspace{1em}
\noindent\emph{JEL Classification:} C14, C58, G21, G28.

\vspace{0.5em}
\noindent\emph{Keywords:} O-minimal structures; Pila--Wilkie
counting; unlikely intersections; systemic risk; cryptocurrency
exchanges; proof of reserves; contagion; natural experiments.

\bigskip

%========================================================================
\section{Introduction}\label{sec:intro}

Cryptocurrency exchanges are among the most transparent
counterparties in modern finance: every one of the top ten publishes
on-chain proof-of-reserves attestations that reveal the composition
of hundreds of billions of dollars of customer assets, and the
underlying wallet balances are independently verifiable at
block-level granularity. Yet no working measure of joint exchange
distress uses this transparency. The paradox is not accidental. It
reflects a fundamental gap in systemic risk measurement: standard
indices are calibrated to institutions that publish audited equity
prices, not to venues that publish audited reserve compositions.

The events of the last three years exposed this gap concretely.
Six major centralised cryptocurrency venues suspended withdrawals,
entered chapter 11, or ceased operations within seven months of
each other: Celsius on June 13, 2022, Voyager Digital on July 5,
2022, FTX and Alameda Research on November 11, 2022, BlockFi on
November 28, 2022, and Genesis Global Capital on January 19,
2023.\footnote{The cluster is documented in the corresponding
chapter~11 filings and independent on-chain forensics of
\citet{Griffin-Kruger-Mei2023}. Voyager and Celsius filings
pre-date the FTX collapse by four to five months, ruling out FTX
as a mechanical trigger for the cluster and motivating the search
for a common systemic factor.} The five principal events left an
aggregate customer-liability shortfall of approximately \$16 billion
(Table~\ref{tab:distress-events}), or roughly $8\%$ of the on-chain
reserves of the three largest surviving venues at the end of 2025.
The cluster is anomalous even by the standards of financial history:
the events are separated by a median of twenty-seven days, share no
common regulator, and predate the March 2023 banking shock by two
months.

Regulators, academics, and market participants attribute the
cluster to contagion, but the operative sense in which the events
cluster remains unclear. Whether the cluster is the tail of a
common systemic factor, a sequence of idiosyncratic failures
amplified by a shared network of counterparties, or the joint
realisation of unlikely intersections in a carefully specified
state space determines both the measure of systemic risk that is
appropriate and the policy that follows from it.

Two edges of the same question follow. What is the sharpest prior on
the frequency of joint distress that public balance-sheet snapshots
can support in the absence of parametric default-probability
assumptions? And what geometric structure of the multi-venue state
space delivers that prior? The regulatory edge and the mathematical
edge, it turns out, admit a common answer. That answer defines a
market-specific ceiling on how rare joint distress ought to be when
no common systemic factor is present---the paper labels this ceiling
the \emph{intersection frontier}, a geometric object derived from a
Pila--Wilkie polylog prior on the reserve-simplex product and distinct
from the connectedness-topological notion of
\citet{DieboldYilmaz2014}---and identifying its location is the
object of the paper.

The paper draws on and contributes to three literatures.
Structural network-contagion models---from \citet{AllenGale2000} and
\citet{FreixasParigiRochet2000}, extended by
\citet{EliottGolubJackson2014}, \citet{AcemogluOzdaglarTahbaz2015},
\citet{GlassermanYoung2015}, \citet{RochetTirole1996},
\citet{CabralesGaleGottardi2016}, and \citet{DuffieZhu2011}---take
observed bilateral exposures as the primitive and study amplification
mechanisms defined on it. They do not deliver a prior on joint
distress when bilateral exposures are unobserved, which is the
generic condition of cryptocurrency venues whose interbank market is
largely off-chain. Systemic-risk measurement---from
\citet{AdrianBrunnermeier2016CoVaR} through
\citet{AchariaEtAl2017SES}, \citet{BrownleesEngle2017SRISK}, and
\citet{BillioEtAl2012}---builds equity-based indices that require
liquid equity data unavailable for most of the top cryptocurrency
exchanges, and are silent on venues that trade only their native
tokens or that operate outside the equity-market perimeter altogether.

Crypto microstructure has documented the on-chain exposures the paper
uses \citep{MakarovSchoar2020, GriffinShams2020,
Griffin-Kruger-Mei2023, FooterKarolyi2023, AlexanderDeng2023,
Cong-Li-Wang2021, FootcaultFuchsGoyalKaplan2024,
LiuTsyvinski2021, GofmanUeda2023,
Aramonte-Huang-Schrimpf2021, Fahlenbrach-Kirchler-Kremer2024},
but has concentrated on price arbitrage rather than joint distress.
\citet{Griffin-Kruger-Mei2023} is the closest neighbour of the
present paper: it uses on-chain forensics to isolate the specific
transaction pattern behind the FTX collapse and delivers an
event-level narrative. The reserve-simplex-product framework of
this paper is complementary in three ways: it produces a market-level
rather than event-level diagnostic, it delivers a directional prior
that a supervisor can evaluate without prosecutorial access to the
failed venue's books, and it is computable in real time on public
proof-of-reserves attestations alone. A parallel supervisory
literature \citep{BCBS2021Crypto, FSB2023Crypto, Aldasoro2023Stablecoin}
frames the same question from the regulator's side but has not had
access to a formal counting theorem on the venues' joint state;
the present paper closes that gap by borrowing from an unlikely
place.

The o-minimality programme in arithmetic geometry
\citep{PilaWilkie2006, vandenDries1998Tame,
Pila2011ManinMumford, BakkerKlinglerTsimerman2020,
PeterzilStarchenko2009}, whose central methodological role in
transferring model-theoretic tools into arithmetic geometry was
cited in Jacob Tsimerman's 2026 Fields Medal work, produced over
the past two decades the sharpest available upper bounds on the
number of rational points of bounded height on the transcendental
part of a definable set. The
related Zilber--Pink family of conjectures
\citep{Pink2005ZilberPink, Zilber2002} organises these bounds under
the heading of \emph{unlikely intersections}; the present paper
deploys the counting theorem itself \citep{PilaWilkie2006}, not any
Zilber--Pink conjecture, and treats the reserve-simplex product as
the finance-native definable set to which the counting theorem is
applied. What is missing from the existing literature is any
statement that ties the counting theorem to the finance-native
reserve-simplex geometry that centralised exchanges disclose.

Closing the gap requires an object that lives natively in both worlds.
The paper constructs it: a reserve-simplex-product representation of
the multi-venue state, built so that (i) every observed proof-of-reserves
attestation is a marked point in the product, (ii) the joint-distress
region becomes a definable subvariety in the o-minimal structure
$\R_{\text{an,exp}}$, and (iii) the Pila--Wilkie counting theorem
delivers a polylog prior on $k$-fold intersections that no
parametric-default model can beat without invoking a common systemic
factor. Above a closed-form crossover threshold in the co-distress
count, the empirical intersection frequency exceeds the prior only if
a common systemic factor is present; below it, isolated failures
suffice. The two regimes meet at an intrinsic $\varepsilon$-loss that
is the necessary cost of scale-induction, inherited directly from the
Pila--Wilkie proof itself.

The paper fills these gaps in three steps that deliver three parallel
contributions.

\emph{First, a reserve-simplex-product representation with a
funding-implied third axis.} A balance-sheet snapshot is
two-dimensional. The paper lifts it to three by constructing a
directional-tube representation of the multi-venue state and
adding a funding-implied axis unique to venues that host perpetual
futures. The intersection frontier lives on this added axis.

\emph{Second, matching o-minimal bounds on the co-distress
frequency.} The paper establishes a pair of sharp, matching bounds
on the frequency of joint distress: from above, a polylog prior no
data-generating process can exceed without a common systemic factor;
from below, a power-of-horizon frequency under any factor of
positive persistence. The two rates meet at an intrinsic
counting-inefficiency loss inherited from the o-minimal upper bound,
the necessary cost that makes the frontier a frontier rather than
a smooth boundary. The crossover threshold is closed-form and
supplies a testable prediction on the regime.

\emph{Third, a counter-intuitive concentration-as-mechanism
finding.} Pooling across venues lifts the empirical intersection
frequency toward the wild regime around a regulatory shock that
raised the safe-asset floor; a stablecoin-issuer placebo shows no
comparable move. Three policy corollaries follow: a real-time
intersection-count audit, a state-contingent capital buffer, and
a coordinated-disclosure recipe.

The remainder of the paper proceeds as follows. Section~\ref{sec:setup}
introduces the reserve-simplex product, the directional-tube
representation, and the o-minimal representation of the joint
distress region. Section~\ref{sec:bounds} proves the bounds on the
intersection frontier. Section~\ref{sec:empirics} carries out the
empirical location of the frontier on the DefiLlama on-chain panel
and the January 2024 spot-Bitcoin ETF natural experiment.
Section~\ref{sec:policy} discusses supervisory, capital-adequacy,
and monitoring implications. Section~\ref{sec:conclusion} concludes.
Proofs and technical lemmas appear in the Appendix. All data and code
are public, archived with SHA-256 hashes, and reproduce every reported
figure, table, and test in the paper (Appendix~\ref{app:data}).

%========================================================================
